\clearpage
\VerdierPhysicalPageAnchor{0221}%
\begingroup\small\setlength{\parskip}{0pt}\setlength{\emergencystretch}{1.5em}
\setlength{\abovedisplayskip}{4pt plus 2pt minus 2pt}%
\setlength{\belowdisplayskip}{4pt plus 2pt minus 2pt}%

\VerdierSemanticAnchor{ch3d-3-3-5}
\medskip\noindent\textbf{Lemme 3.3.5.}\quad
\textit{Soient $m\colon W\longrightarrow W'$, $m'\colon W'\longrightarrow W''$,
$m''=m'm$ trois morphismes de $\mathcal A^{\operatorname{Ext}}$ appartenant
à $\mathcal F\ell_{\leq 1}(\mathcal A^{\operatorname{Ext}})$. Les morphismes
$F(m)$, $F(m')$ et $F(m'')$ ont été définis ci-dessus dans les cas 1), 2),
3), 4), 5), 6) et on a :}

\VerdierSemanticAnchor{ch3d-3-3-5-1}
\begin{equation}
F(m'')=F(m')F(m).
\tag{3.3.5.1}
\end{equation}

\VerdierSemanticAnchor{ch3d-3-3-6}
\medskip\noindent\textbf{Proposition 3.3.6.}\quad
(Définition de la catégorie $\mathcal A^{\operatorname{Ext}}$ par générateurs
et relations). \textit{Soient $\mathcal C$ une catégorie et :}

\VerdierSemanticAnchor{ch3d-3-3-6-object-map}
\[
\operatorname{Ob}(F)\colon
\operatorname{Ob}(\mathcal A^{\operatorname{Ext}})
\longrightarrow \operatorname{Ob}(\mathcal C),
\]

\VerdierSemanticAnchor{ch3d-3-3-6-arrow-map}
\[
\mathcal F\ell_{\leq 1}(F)\colon
\mathcal F\ell_{\leq 1}(\mathcal A^{\operatorname{Ext}})
\longrightarrow \mathcal F\ell(\mathcal C)
\]
\textit{deux applications compatibles avec les applications source et but
des catégories $\mathcal A^{\operatorname{Ext}}$ et $\mathcal C$ telles que,
pour tout triplet de morphismes $m\colon W\longrightarrow W'$,
$m'\colon W'\longrightarrow W''$, $m''=m'm$,
$m,m',m''\in\mathcal F\ell_{\leq 1}(\mathcal A^{\operatorname{Ext}})$,
on ait :}

\VerdierSemanticAnchor{ch3d-3-3-6-relation}
\[
\mathcal F\ell_{\leq 1}(F)(m')\,
\mathcal F\ell_{\leq 1}(F)(m)
=\mathcal F\ell_{\leq 1}(F)(m'').
\]

\VerdierSemanticAnchor{ch3d-3-3-6-extension}
\textit{Il existe alors un foncteur
$F\colon\mathcal A^{\operatorname{Ext}}\longrightarrow\mathcal C$ et un seul
tel que $F$ induise l'application $\operatorname{Ob}(F)$ sur les ensembles
d'objets et tel que, pour tout
$m\in\mathcal F\ell_{\leq 1}(\mathcal A^{\operatorname{Ext}})$, on ait :}
\[
F(m)=\mathcal F\ell_{\leq 1}(F)(m).
\]

En effet, on vérifie facilement que la suite de foncteurs liés associée au
foncteur fourni par (3.3.5) et (3.3.6) est bien la suite
$(F^n,\delta_F^n)$ qu'on avait au départ.

\VerdierSemanticAnchor{ch3d-3-3-7}
\medskip\noindent\textbf{3.3.7. Démonstration du lemme 3.3.5.}\quad
La démonstration se fait en examinant les différents cas possibles. Il y a,
d'abord, quatre cas à examiner lorsqu'au moins deux des objets $W$, $W'$,
$W''$ sont nuls. Les morphismes $F(m)$, $F(m')$ et $F(m'')$ sont alors
définis dans l'un des cas 1), 2) ou 3). La vérification de (3.3.5.1) est
immédiate et laissée au lecteur. Il y a, ensuite, neuf cas à examiner dans
lesquels un des objets $W$, $W'$, $W''$ est nul. Deux des morphismes
$m$, $m'$, $m''$ sont alors définis par 1), 2), 3), le troisième étant
défini par 4), 5), 6). Là encore la vérification est immédiate et laissée au
lecteur. Il reste à examiner les cas où $W$, $W'$, $W''$ sont non nuls.
Ces cas sont au nombre de huit :

\VerdierSemanticAnchor{ch3d-3-3-7-cases-start}
\noindent$c_1)$ degré de $m$ et $m'<0$ ;

\noindent$c_2)$ degré de $m<0$, degré de $m'=0$ ;
\endgroup

